\documentclass[12pt,letterpaper]{article}
\title{A Walking Tour of the Old Harbour}
\author{Cornelius Whitlock}
\date{August 20, 2021}
\begin{document}
\maketitle

\section{Beginning at the Quay}
The tour starts where the ferries once docked, beside a weathered
lighthouse that no longer guides any ship. Visitors are advised to wear
comfortable shoes, since the cobblestones are uneven and frequently wet.

\begin{quote}
The harbour is not a place but a conversation between the land and the sea,
and neither side has ever finished speaking.
\end{quote}

\subsection{The Fish Market}
Early risers will find the market at its liveliest. Traders call out prices
in a rapid, incomprehensible dialect, and the unwary purchaser may discover
that haggling is expected.

\subsection{The Customs House}
The handsome building opposite was constructed in the nineteenth century.
Its facade, though blackened by decades of coal smoke, retains remarkable
ornamental detail.\footnote{Guided tours of the interior run on Saturdays.}

\section{Along the Breakwater}
A long walk follows, rewarded by an unobstructed view across the bay.

\begin{quotation}
Local tradition holds that anyone who walks the breakwater at dusk and
counts the fishing boats will always arrive at an odd number. No one has
satisfactorily explained this, though many have tried.

The author counted eleven.
\end{quotation}

\section*{Practical Notes}
\begin{description}
\item[Opening hours] The museum opens at ten and closes at five.
\item[Refreshments] A small caf\'e near the lighthouse serves soup and bread.
\item[Weather] Bring a waterproof jacket regardless of the forecast.
\end{description}

\vspace{1cm}
\begin{center}
\emph{Fair winds to all who wander here.}
\end{center}
\end{document}
